% !TEX root = ../main.tex
% =====================================================================
% 6. DISCUSSION
%
% Section 5 says WHAT happened; this section says WHAT IT MEANS.
% Only conclusions the tables and figures actually support, and no
% repetition of the results themselves.
% =====================================================================

\section{Discussion}
\label{sec:discussion}

% ---------------------------------------------------------------------
\subsection{Why the parameter reduction did not cost accuracy}
\label{sec:disc:worked}

The bottleneck encoder removes parameters from one specific place, and the results indicate it was
the right place. After a convolution is split into a depthwise and a pointwise step, the pointwise
step dominates: at $\dmodel = 128$ it costs three times as much as all three depthwise convolutions
combined (\Cref{sec:arch:block}). The depthwise step looks along time; the pointwise step only mixes
channels. \Cref{eq:bottleneck} halves the second and leaves the first untouched, which is why the
encoder can lose 28\,\% of its weights without losing the mechanism that extracts temporal patterns.
What the results add is that most of the channel-mixing capacity was not being used productively:
\mbot{} matches or exceeds the re-trained \millet{} baseline on accuracy, AOPCR and NDCG@$n$
simultaneously with 90\,\% fewer parameters (\Cref{tab:web_main}). That is a claim about the
baseline's capacity as much as about the proposed encoder, and the parameter arithmetic alone could
not have predicted it.

% ---------------------------------------------------------------------
\subsection{Why Top-$k$ pooling improved the explanations}
\label{sec:disc:topk}

Top-$k$ pooling produced the best NDCG@$n$ of every model with complete seed data, and the mechanism
can be identified rather than inferred. Averaging only the $\kappa \nsteps$ largest evidence values
makes $\interp_{k,j}$ exactly zero everywhere the model did not use, instead of a small weight that
never quite reaches zero; since NDCG@$n$ scores whether the truly
modified steps appear at the top of the ranking, removing the background from that ranking is
directly what the metric rewards.

The $\kappa$ ablation (\Cref{tab:ablation_topk}) confirms this and bounds it. NDCG@$n$ is highest at
$\kappa = 0.1$ (0.777) and falls monotonically as $\kappa$ grows towards 1 (0.722, the value of
class-wise conjunctive pooling), exactly the behaviour the design predicts. Accuracy is not
monotone: it peaks between $\kappa = 0.1$ and $0.25$ and collapses to 0.888 at $\kappa = 0.05$, so
selecting too few time steps discards evidence the classifier needs and the mechanism has an optimum
rather than improving without limit. The complementary effect appears in the attention entropy term:
switching it off raises accuracy from 0.938 to 0.950 while lowering AOPCR from 2.778 to 2.303 and
NDCG@$n$ from 0.777 to 0.765 (\Cref{tab:ablation_focus}). Both experiments point the same way --
concentrating the evidence improves the explanation and costs about one accuracy point.

% ---------------------------------------------------------------------
\subsection{Encoder and pooling head interact}
\label{sec:disc:interaction}

The best model of the study, \mgat{}, is not built from the individually best encoder and the
individually best pooling head: averaged over their pairings, the gated encoder reaches 0.902 and
Top-$k$ pooling 0.904, both mid-table (\Cref{tab:ablation_encoder}, \Cref{tab:ablation_pooling}).
Two ordinary components producing the best combination is evidence of an interaction. A plausible
mechanism is that the gated encoder's channel gate (\Cref{eq:gate}) already suppresses the channels
carrying no signal, making a hard-selection head the natural partner, whereas a smoother head would
reintroduce the neighbouring evidence the gate has just attenuated. This is consistent with the grid
but not established by it: the one-factor-at-a-time experiment that would separate this mechanism
from alternatives was not run. The practical consequence is that encoder and pooling head cannot be
optimised independently, and since neither component's spread dominates the other's
(\Cref{sec:res:ucr}), the combination has to be searched rather than composed.

% ---------------------------------------------------------------------
\subsection{Where the approach failed}
\label{sec:disc:negative}

\paragraph{Dual-stream pooling: a safety property that training never exercises.}
This is the clearest negative result. Property~\ref{prop:ds} states that dual-stream pooling cannot
be worse than the conjunctive baseline, because its blend $\lambda$ can reach 0 and recover it
exactly. In practice it is the weakest head tested by a wide margin: mean \webtraffic{} accuracy
0.744 over its four encoder pairings against 0.90 and above for class-wise, softmax conjunctive and
Top-$k$ pooling, with one pairing collapsing to 0.556 (\Cref{tab:ablation_pooling}).

The property is about representational capacity, while the failure is one of optimisation. No term
in \Cref{eq:loss} pushes $\lambda$ towards 0, so if the critical-instance stream of
\Cref{eq:ds_query} learns a poor query early in training, gradient descent has no incentive to return
to the safe mean. The model \emph{can} represent the baseline; it never finds it.

\paragraph{Per-class gated attention: no fallback at all.}
Gated attention is a weaker version of the same problem: it is the only proposed head with no blend
parameter, so no setting of its weights reduces it to conjunctive pooling (\Cref{sec:pool:gated}).
Its mean \webtraffic{} accuracy of 0.918 is respectable, but it is never the best partner for any
encoder, unlike Top-$k$ and class-wise pooling which each produce a top-ten model. Without a
known-good configuration to fall back on, a poorly trained gate has nothing to recover to, and the
head settles in the middle of the ranking rather than either winning or failing outright.

\paragraph{What both cases indicate.}
A property that holds by construction is only useful if training actually puts the model in the
state where it applies. Neither head got there under the single configuration used across the whole
study (\Cref{sec:setup:protocol}). Two fixes follow directly: start $\lambda$ near 0 so the model has
to move away from the safe branch on purpose, or add a penalty that holds a blended head close to
that branch early in training.

% ---------------------------------------------------------------------
\subsection{How the numbers should be read}
\label{sec:disc:reading}

\paragraph{AOPCR compares rows of one table, and nothing else.}
AOPCR is an area under a curve of raw confidence differences, so its scale follows the scale of the
model's logits and has no fixed range (\Cref{sec:bg:metrics}). Every model in \Cref{tab:web_main}
was trained under one configuration, so its AOPCR column is internally comparable; an AOPCR produced
under a different configuration is not on the same scale and must not be placed beside it.
\Cref{app:published} measures how large that effect is on one architecture, and it is large.

\paragraph{Differences smaller than the seed spread are not results.}
Where three seeds are available the spread is not negligible: accuracy standard deviation ranges
from 0.003 for \mgat{} to 0.030 for \minp{}, and AOPCR is proportionally noisier still -- 0.437 for
\minp{}, about a sixth of its own mean, and 0.884 for the re-trained \millet{} baseline
(\Cref{tab:web_main}). This was applied as a rule in both directions. \minp{}'s AOPCR advantage over
\mbot{} (2.651 against 2.621) is far smaller than either model's standard deviation and is reported
as noise, not as a result; equally, \mgat{}'s accuracy lead over \mbot{} (0.955 against 0.947) cannot
be called decisive, because the second model's spread is 0.016. For the single-seed models no such
test is possible, and those numbers are indicative rather than exact.

\paragraph{The training budget was the same for everyone, and it was short.}
The configuration of \Cref{tab:protocol} trains for at most 400 epochs with patience 60. That was
the budget 72 variants over 129 datasets allowed (\Cref{app:constraints}), and it applies equally to
the baselines and to the \seanet{} models, so it does not favour either. It does mean the absolute
accuracies here are lower than what a longer schedule would produce, on both sides.
\Cref{app:published} quantifies that on one architecture, and re-training the \seanet{} encoders
under a longer schedule is the first item of future work.

% ---------------------------------------------------------------------
\subsection{Limitations}
\label{sec:disc:limitations}

Six limitations bound what the results above support. \textbf{Explanation correctness is measured on
one dataset}: NDCG@$n$ needs per-time-step ground truth, which only \webtraffic{} provides, and
\webtraffic{} is synthetic, so its important regions are generated rather than observed and it is
not known how the result holds on real, noisier data. \textbf{The headline models were selected on
that same dataset}, and \Cref{sec:res:ucr} shows what that costs: over the 85 \ucr{} datasets the
re-trained baseline is ahead of both. \textbf{Most models have a single seed} -- four have 3 seeds
on \webtraffic{} and two of those also over the \ucr{} archive, while the other 68 have one; where
repeats exist the spread was large enough to change how a claim is phrased
(\Cref{sec:disc:reading}), which is a reason for caution about the single-seed numbers rather than a
reason to trust them. \textbf{Only univariate series were tested}: the encoder is indifferent to the
number of input channels, so multivariate input should work mechanically, but that was not verified.
\textbf{Faithfulness is not usefulness}, because AOPCR and NDCG@$n$ both measure whether the stated
explanation matches what the model used, and neither establishes that a person would find the
highlighted time steps informative -- that needs a human study. Finally, \textbf{nothing has run on
a real microcontroller}: the cost measurements indicate that it should fit (\Cref{tab:cost}), but
that is not the same as a working deployment (objective~O6, \Cref{tab:objectives}).
